\chapter{Panoramica di $\pi$Tantum}
\label{ch:panoramica_pitantum}

\epigraph{``A program does what you tell it to do, not what you
want it to do.''}{Maxim popolare attribuita a vari}

\lettrine[lines=3]{P}{}rima di entrare nei dettagli di
ciascun metodo \`e utile guardare $\pi$Tantum dall'alto: che
cosa fa, com'\`e organizzato, da quali pezzi \`e composto, e
con quale flusso un coordinatore d'orario lo usa per
costruire l'orario annuale di una scuola.

\section{Cosa fa $\pi$Tantum (in tre frasi)}

$\pi$Tantum \`e un sistema software che, ricevuti in input
l'anagrafica di una scuola (docenti, classi, materie, aule,
indirizzi, gruppi, studenti), i monte ore per indirizzo, le
preferenze e le indisponibilit\`a, produce in output un
\emph{orario settimanale completo} che soddisfa tutti i
vincoli rigidi (HARD) e minimizza la somma pesata delle
violazioni dei vincoli morbidi (SOFT). Una volta prodotto
l'orario, lo stesso sistema lo gestisce in corso d'anno:
modifiche manuali con controllo di fattibilit\`a in tempo
reale, gestione delle assenze e delle supplenze quotidiane,
visualizzazione per docente / per classe / per aula, export
in Excel/PDF.

\section{Tre layer in un unico repository}

\sidenote{Il repository \texttt{gborghi/timetable} \`e
organizzato in tre directory principali, una per ogni layer.}

Il sistema \`e organizzato in tre strati distinti, ciascuno
con il proprio dominio di responsabilit\`a.

\subsection{Solver layer (\texttt{experiments/})}

Codice Python puro che implementa la pipeline di
ottimizzazione. Non sa nulla di database, n\'e di interfacce
utente: legge un dizionario di input (docenti, classi,
vincoli) e produce un dizionario di output (orario). Le
componenti principali sono:
\begin{itemize}
\item \texttt{cpsat\_v2\_assignment.py} (Phase A): matching
  docente $\to$ classe $\to$ ore.
\item \texttt{cpsat\_v2\_timetable.py} (Phase B): timetabling
  CP-SAT su singolo cluster.
\item \texttt{decomposition\_spectral\_v2.py}: clustering
  spettrale del grafo bipartito classe--docente.
\item \texttt{metaheuristics.py}: LNS, SA, TS, ILS che
  raffinano la soluzione di Phase B.
\item \texttt{alns.py}, \texttt{vns.py}: varianti adattive e
  a vicinato variabile.
\item \texttt{hall\_check.py}, \texttt{lagrangian.py},
  \texttt{column\_generation.py}: tecniche aggiuntive.
\item \texttt{classroom\_assignment.py} (Phase D): assegnazione
  delle aule alle (classe, slot).
\end{itemize}

\subsection{Web UI layer (\texttt{webui/})}

Backend FastAPI + frontend SvelteKit. Espone tutte le
operazioni del solver come API REST e fornisce un'interfaccia
web a sedici tab che copre l'intero ciclo di vita di un
orario: dall'inserimento dei dati anagrafici (Docenti,
Classi, Aule, Materie, Indirizzi, Studenti, Gruppi) alla
configurazione dei vincoli (matrici 5-stati, vincoli logici
DSL), alla generazione (Workflow), all'editing (Orario), alla
gestione operativa (Assenze e supplenze, Monitor).

Il frontend e il backend comunicano via HTTP/JSON. Lo storage
\`e SQLite (un singolo file
\texttt{webui/data/timetable.db}); le migrazioni sono
idempotenti, eseguite a ogni avvio del backend.

\subsection{Legacy layer (\texttt{schedule/})}

Script monolitici che hanno fatto da prototipo per il solver
attuale. Sono mantenuti per riferimento storico e per
benchmark riproducibili contro le versioni precedenti, ma non
sono pi\`u parte del flusso operativo.

\section{Diagramma a blocchi}

\begin{figure}[h]
\centering
\begin{tikzpicture}[node distance=1.0cm and 0.8cm,
  block/.style={draw, very thick, rounded corners=4pt,
                minimum width=4cm, minimum height=1.0cm,
                align=center, font=\small\sffamily,
                fill=cBox!40},
  arr/.style={-{Stealth[length=2.5mm]}, thick}]
  \node[block] (fe) {Frontend SvelteKit + Vite\\\footnotesize 16 tab, drag-and-drop};
  \node[block, below=of fe] (be) {Backend FastAPI\\\footnotesize 18 router, $\sim$118 endpoint};
  \node[block, below=of be] (db) {SQLite\\\footnotesize $\sim$30 tabelle};
  \node[block, fill=cOro!25, below=of db] (en) {Solver layer (Python)\\\footnotesize CP-SAT + spectral + meta};
  \node[block, fill=cIndaco!15, right=2.0cm of en] (or) {OR-Tools (CP-SAT)};
  \draw[arr] (fe) -- node[draw=none, right] {\scriptsize \texttt{/api/*}} (be);
  \draw[arr] (be) -- node[draw=none, right] {\scriptsize SQLAlchemy} (db);
  \draw[arr] (be) to[bend left=45]
        node[draw=none, right] {\scriptsize \texttt{engine\_io.py} (pickle)} (en);
  \draw[arr] (en) -- (or);
\end{tikzpicture}
\caption{Architettura a blocchi di $\pi$Tantum: il frontend
parla al backend via REST, il backend persiste su SQLite e
chiama il solver via pickle, il solver invoca OR-Tools per
CP-SAT.}
\label{fig:arch_panoramica}
\end{figure}

\section{Flusso d'uso tipico}

Un coordinatore d'orario, alla prima esecuzione, segue
tipicamente questo flusso:

\begin{enumerate}
\item \textbf{Importa l'anagrafica}: pu\`o partire da template
  Excel/CSV (un foglio per entit\`a: docenti, classi, materie,
  aule, indirizzi, studenti, gruppi). Oppure usa la dashboard
  ``Importa profilo'' per partire da un dataset mock di
  taglia \texttt{small}/\texttt{medium}/\texttt{big}.
\item \textbf{Verifica e completa i dati}: tab Docenti,
  Classi, etc., per controllare che tutto sia coerente. Le
  matrici 5-stati permettono di esprimere indisponibilit\`a
  (HARD) e preferenze (SOFT/PREF) per ogni cella
  (giorno, ora). Le matrici delle aule vincolano quali
  materie possono andare in quali aule.
\item \textbf{Aggiunge vincoli logici}: il DSL generico
  (cap.~\ref{ch:metodo_dsl}) permette di esprimere vincoli
  su qualunque combinazione di entit\`a. Esempio:
  ``\texttt{forall t in teachers where t.subject == "Fisica":
  count l in lessons where l.teacher == t and
  l.classroom.type == "lab\_fisica": l == 1}''.
\item \textbf{Lancia il pre-check di Hall}: bottone in tre
  punti dell'UI. Se Hall fallisce, il sistema mostra il
  sotto-insieme di classi che ha pi\`u domanda che offerta;
  l'utente capisce \emph{quale} vincolo va rilassato.
\item \textbf{Lancia la pipeline}: tab Workflow, bottone
  ``Pipeline completa''. Phase A (assignment) + Phase B
  (decomposizione spettrale + CP-SAT) + Phase C (LNS+SA+TS+ILS)
  + Phase D (aule). Output: orario completo, log dettagliato,
  telemetria runtime.
\item \textbf{Editor manuale}: tab Orario. Drag-and-drop di
  lezioni; il sistema controlla in tempo reale che la mossa
  sia HARD-fattibile e mostra come cambia la funzione
  obiettivo (delta SOFT).
\item \textbf{Gestione corso d'anno}: tab Assenze e
  supplenze. Marca i docenti assenti, vede le classi
  scoperte, drag-drop dei docenti disponibili come
  supplenti.
\end{enumerate}

\begin{casostudiobox}
\textbf{Liceo ``Galileo'', 32 classi, settembre}. Il
vicepreside parte da un import Excel con 67 docenti,
32 classi, 11 indirizzi (4 scientifico, 3 classico, 2
linguistico, 2 scienze applicate). Aggiunge 4 vincoli
logici (es. ``il prof Bianchi non lavora sabato'', ``il lab
di fisica solo nel triennio'') e lancia la pipeline. In
9~minuti $\pi$Tantum produce un orario completo che
soddisfa tutti i vincoli HARD e minimizza una funzione
obiettivo composta da 7 termini SOFT pesati.

Un'occhiata al log mostra che Phase~A ha impiegato 12
secondi (assignment trovato facilmente), Phase~B 4~minuti
(la scuola \`e stata spezzata in 5 cluster da 6--8 classi
ciascuno), Phase~C 3~minuti di LNS + 1~minuto di SA, Phase~D
40~secondi (aule). L'obiettivo finale \`e a 11.2 punti, con
0 violazioni HARD e 23 violazioni SOFT (di cui 14 sono ``ore
isolate'' che il vicepreside decide di accettare).
\end{casostudiobox}

\section{Quando $\pi$Tantum non basta da solo}

\begin{erroricomunibox}
\begin{itemize}
\item \textbf{Dati di input incompleti o contraddittori}.
  Se un docente \`e abilitato ad A027 ma il suo monte ore
  supera la somma delle ore di matematica disponibili, il
  pre-check di Hall fallisce. \`E un \emph{problema dei dati},
  non del solver: $\pi$Tantum non pu\`o ``inventare'' classi.
\item \textbf{Vincoli sociali impliciti}. ``La prof Rossi
  vorrebbe avere tutte le sue ore al mattino'' richiede
  \emph{esplicitazione}: matrice di preferenza con celle
  pomeridiane in SOFT. Se la richiesta non \`e codificata,
  il solver non la conosce.
\item \textbf{Iterazione mancante con il committente}. Una
  prima soluzione produce sempre richieste di modifica.
  $\pi$Tantum ha un editor manuale e un sistema di vincoli
  logici incrementale apposta: la pipeline non \`e da
  rilanciare ogni volta da zero.
\end{itemize}
\end{erroricomunibox}

\section{Tecnologie scelte (e perch\'e)}

\sidenote{Tutte le tecnologie scelte sono \emph{open source}
e prive di lock-in commerciali. SQLite, FastAPI, SvelteKit,
OR-Tools sono progetti maturi e attivamente mantenuti.}

\begin{itemize}
\item \textbf{Python 3.11+}: linguaggio dominante per il
  scientific computing e per la modellazione constraint.
  Performance accettabili su questo dominio (la parte
  hot \`e il solver C++ di OR-Tools, non il codice Python
  che lo guida).
\item \textbf{OR-Tools} (Google): solver CP-SAT
  \emph{state-of-the-art}, gratuito, con API Python pulite.
\item \textbf{FastAPI}: backend async-friendly, validazione
  automatica via Pydantic, documentazione OpenAPI/Swagger
  generata automaticamente.
\item \textbf{SQLite}: zero-config, single-file, perfetto
  per un sistema mono-utente o piccolo-team.
\item \textbf{SvelteKit + Tailwind}: stack frontend leggero,
  hot-reload veloce, bundle finale piccolo. Svelte~5 in
  legacy mode.
\item \textbf{Vite}: dev server velocissimo, proxy nativo
  verso il backend.
\end{itemize}

Le alternative considerate (Django, Flask, React, Next.js,
Gurobi commerciale, MiniZinc) sono state scartate per
ragioni di leggerezza, costo (zero), facilit\`a di
deployment in scuole italiane (un solo file binario, niente
licenze).

\section{Roadmap del libro}

I prossimi capitoli ripercorrono in dettaglio gli aspetti
toccati qui in panoramica.

La parte~II descrive le entit\`a (cap.~\ref{ch:vincoli}, cap.~\ref{ch:guida_ui})
e i casi d'uso reali. La parte~III \`e il cuore tecnico:
CP-SAT (cap.~\ref{ch:metodo_cpsat}), decomposizione spettrale
(cap.~\ref{ch:metodo_spettrale}), metaeuristiche
(cap.~\ref{ch:metodo_metaeuristiche}), Hall pre-check
(cap.~\ref{ch:metodo_hall}), Lagrangian + Column Generation
(cap.~\ref{ch:metodo_lagrangian}), Monte Carlo
(cap.~\ref{ch:metodo_montecarlo}), analisi del grafo
bipartito (cap.~\ref{ch:metodo_grafo}), statistica applicata
(cap.~\ref{ch:metodo_statistica}), DSL
(cap.~\ref{ch:metodo_dsl}). La parte~IV documenta
l'architettura, il modello dati, le API REST, l'estensione.
La parte~V riporta i benchmark sui cinque profili e una
sezione di \emph{lessons learned}.

\pullquote{Un sistema di scheduling \`e tanto buono quanto la
sua capacit\`a di accettare modifiche dopo la prima
esecuzione.}{prassi consolidata fra i coordinatori d'orario}
